% Z-005 — Resposta do assistente (SPEC-009, REQ-AI-001..005)
%
% VERIFICAÇÃO: parse + typecheck com CZT via scripts/verify-z.ps1. NÃO prova teoremas;
% as conjecturas (\vdash?) são checadas por Hypothesis em tests/unit/assistant/test_invariants.py.
\documentclass{article}
\begin{document}

\section*{Z-005: Assistant answer}

\begin{zed}
  [CHUNK, DOC, EQTYPE, EXCERPT]
\end{zed}

\begin{zed}
  OUTCOME ::= answered | insufficientEvidence | refused \\
  RISK ::= normalRisk | highRisk \\
  REVIEW ::= reviewRequired | reviewNotRequired
\end{zed}

\begin{schema}{KnowledgeBase}
  approved : \finset DOC \\
  chunkDoc : CHUNK \pfun DOC \\
  chunkRisk : CHUNK \fun RISK \\
  scope : CHUNK \fun \finset EQTYPE \\
  text : CHUNK \fun \finset EXCERPT
\where
  \ran chunkDoc \subseteq approved
\end{schema}

\begin{schema}{Answer}
  KnowledgeBase \\
  equipmentType? : EQTYPE \\
  retrieved : \finset CHUNK \\
  outcome! : OUTCOME \\
  sources! : \finset CHUNK \\
  excerpts! : \finset EXCERPT \\
  review! : REVIEW
\where
  retrieved \subseteq \dom chunkDoc \\
  \forall c : retrieved @ equipmentType? \in scope~c \\
  sources! \subseteq retrieved \\
  excerpts! \subseteq \bigcup \{ c : sources! @ text~c \} \\
  outcome! = answered \iff sources! \neq \emptyset \\
  outcome! \neq answered \implies excerpts! = \emptyset \\
  outcome! \neq answered \implies review! = reviewRequired \\
  (\exists c : sources! @ chunkRisk~c = highRisk) \implies review! = reviewRequired
\end{schema}

\subsection*{Conjecturas}

A1: toda fonte citada pertence a um documento aprovado.

\begin{zed}
  \vdash? \forall Answer @ \forall c : sources! @ chunkDoc~c \in approved
\end{zed}

A2: toda fonte citada está no escopo do tipo de equipamento consultado.

\begin{zed}
  \vdash? \forall Answer @ \forall c : sources! @ equipmentType? \in scope~c
\end{zed}

A3: resposta sem fonte nunca é \emph{answered} (sem evidência não há orientação).

\begin{zed}
  \vdash? \forall Answer @ sources! = \emptyset \implies outcome! \neq answered
\end{zed}

% A4 (recusa nunca cita fonte) segue do esquema: outcome! \neq answered => excerpts! = \emptyset.
% A5 (cada excerto é literal) é a propriedade \subseteq \bigcup acima; o verificador de texto em
% spatial_assistant.grounding aplica o análogo lexical à saída de um LLM.

\end{document}
